\section{Discussion}\label{sec:discussion}

\subsection{Managerial implications}

\paragraph{Protect scarce clean capacity} The largest operational loss in multi-site carbon-aware scheduling does not come from imperfect carbon forecasts. It comes from spending low-carbon capacity on work that could have waited or run elsewhere. Proposition~\ref{prop:myopic} and Theorem~\ref{thm:lb} show this in the worst case, and the 2024 experiments confirm it empirically: with migration allowed, a myopic planner with perfect foresight of carbon intensity is still 16--20\% above the optimum. For operators, clean capacity at a low-carbon site should be treated as a scarce resource with an opportunity cost, much as yield managers protect seats for high-fare late bookings. CARMA sets this protection level implicitly, through ghost jobs learned from the operator's own history, and needs no hand-tuned reservation thresholds.

\paragraph{Invest in forecasts where they pay} The value of better intensity forecasts depends on the rest of the policy. For a myopic multi-site planner, perfect forecasts are worth only 1--3 percentage points. Once demand is anticipated, they are worth 4--7 points. When work cannot migrate, they are worth almost the whole remaining gap. Forecasting effort, for example the integration of numerical weather predictions~\cite{maji2022carboncast,lowry2018day}, should therefore follow, not precede, the fix of demand myopia in multi-site settings. In single-site temporal shifting it is the main lever.

\paragraph{Use risk premia with care} Calibrated uncertainty is valuable, and ACI kept coverage on target through a year of strong decarbonisation drift. Its best use in scheduling is not obvious, however. A pessimistic carbon price improved a myopic planner by acting as a crude substitute for anticipation, but it degraded CARMA and hurt temporal-only scheduling, where it suppressed good deferrals. Because the planner re-optimises every hour, much of the protection that a risk premium would buy is already provided by recourse. Calibrated quantiles are better suited to reporting and to service-level decisions than to inflating costs by default.

\paragraph{Guarantees an operator can rely on} Worst-case properties matter for adoption, because an operator must bound the damage that a wrong forecast or an unusual week can do. Theorem~\ref{thm:carma} gives three such bounds. With a backstop resource, bad predictions can never make a job late. Emissions are never more than a factor $\theta$ above the clairvoyant optimum. Excess emissions grow at most linearly with prediction error, and the controlled-error experiments confirm this degradation profile. Theorem~\ref{thm:lb} also sets expectations: no scheduler can promise near-optimal worst-case emissions without some knowledge of future demand, so maintaining a demand profile is a necessity, not a refinement.

\paragraph{Fair carbon accounting for shared capacity} Pooling halved the emissions of the four sites, but location-based accounting, which charges each site for its own jobs wherever they run, would have made South Scotland and North West England worse off than on their own in almost every week. A multi-organisation or multi-tenant pool attributed this way is unstable, because its cleanest members have an incentive to leave. The dual attribution of Theorem~\ref{thm:core} charges each site the carbon shadow price of its demand and credits it for the capacity it contributes. It is budget-balanced, no coalition would prefer to secede, and it is a by-product of the scheduling problem. It offers a principled basis for internal carbon pricing, for sharing the benefits of carbon-aware scheduling among business units, and for tenant-level carbon reporting in shared data centres.

\paragraph{Deployment} CARMA needs three inputs: a public carbon-intensity feed, a demand profile computed from the operator's own job logs, and a network-flow solver. It decides within tens of milliseconds and returns a plan that is optimal for its model, which makes its decisions auditable. This supports the transparency expected of carbon claims in sustainability reporting. The algorithm is also a rare case in which lower emissions come with better service: under heavy load, CARMA finished far more work on time than the other forecast-driven planners (0.05\% late work, against 1.2--2.0\%).

\subsection{Limitations and future research}

The study has several limitations. First, workloads are synthetic. The generator has realistic features (diurnal and weekly seasonality, heavy-tailed job sizes, a mix of urgent and flexible jobs, width limits), but results on production traces may differ. In particular, CARMA's advantage depends on how predictable arrivals are, and the sensitivity analysis shows how it degrades when the load deviates from the learned profile. Second, we use average (attributional) regional intensities as published by NESO. Marginal emission factors may be more appropriate for evaluating load shifting~\cite{dodge2022measuring}, and regional values in Great Britain are modelled estimates, not metered data. Third, the facility parameters (PUE, network energy) are stylised, idle server power is not modelled, and embodied emissions are outside the scope of this study~\cite{acun2023carbon}. Fourth, the forecaster deliberately excludes weather forecasts, so its day-ahead skill is limited. Fifth, bag-of-tasks jobs with free preemption and migration are an idealisation. Jobs with state, data-gravity or data-sovereignty constraints would add restrictions that break the pure network structure. Sixth, the guarantees of Theorem~\ref{thm:carma} assume a backstop resource and a positive emission floor, and the smoothness bound assumes that the planning window covers the episode. The simulations use neither a backstop nor episodes that short, so they test the predicted behaviour rather than the bound itself. The core property of Theorem~\ref{thm:core} is exact for the hindsight benchmark; for CARMA's realised emissions it held in all but 5 of 728 coalition checks.

Several extensions follow naturally. The ghost-demand profile could be replaced by probabilistic arrival forecasts, with the anticipation weight set by a newsvendor-type argument. Scenario-based or distributionally robust variants could treat demand and carbon uncertainty jointly~\cite{bertsimas2004price}. The framework can also be coupled with on-site renewables and storage~\cite{souza2023ecovisor,acun2023carbon} and with demand-response programmes~\cite{wierman2014opportunities}. Finally, the robustness guarantee could be tightened from $\theta$ towards the $\Theta(\sqrt\theta)$ lower bound, for example by protecting deferral decisions with reservation prices~\cite{elyaniv2001optimal,lechowicz2025learning} while retaining consistency, and the attribution could be extended to online, per-decision carbon prices.

\section{Conclusions}\label{sec:conclusion}

This paper studied carbon-aware scheduling of deferrable computing workloads across data-centre sites as a sustainable-operations problem. We showed that the offline problem, including width limits, deadlines and the emissions of data transfer, is a minimum-cost flow. This yields exact oracles and fast online planning. We proved that uncertainty about demand alone forces every online algorithm to a competitive ratio of order $\sqrt\theta$, even with perfect carbon forecasts, and that myopic planning does far worse. We proposed CARMA, which augments the planning network with ghost jobs for expected arrivals, and a probabilistic carbon-intensity forecaster with adaptive conformal calibration. We proved that CARMA meets every deadline and is $\theta$-competitive under any predictions, optimal under exact predictions and smooth in between. The dual prices of the pooled problem also yield a carbon attribution that is in the core of the pooling game. In 52 out-of-sample weeks of 2024 on British regional grid data, CARMA cut emissions relative to carbon-agnostic operation by 36--68\% when migration was allowed. That is within 6--9\% of the clairvoyant optimum and 7.6--9.4\% below myopic MPC, and CARMA also reduced late work. Anticipating demand was worth far more than perfect carbon foresight, while risk-averse carbon pricing helped only myopic planners. CARMA degraded gracefully under injected prediction errors. Location-based carbon accounting violated coalitional stability in every test week. The dual attribution never did for the hindsight benchmark, and failed only 5 of 728 checks for CARMA's realised emissions. When jobs cannot migrate, achievable reductions are smaller (9--18\%) and forecast accuracy is the binding constraint. Carbon-aware computing is most effective when schedulers plan for the work that has not yet arrived, not only for the grid conditions that lie ahead.
